\documentclass{scrartcl}
\usepackage{style}

\newcommand{\versionmajor}{0}
\newcommand{\versionminor}{1}
\newcommand{\versionpatch}{0}
\newcommand{\version}{\versionmajor.\versionminor.\versionpatch}

\title{\LARGE
    Language Model-Based Framework for Automated Data Extraction and Quality Control in Oncological Cancer Registries
}

\subtitle{Ph.D. Project Proposal}

\author{
    Marco Raggini
}

\date{
    \today
}

\begin{document}

\maketitle

\paragraph{Today's oncology research}{In recent years, oncology research has increasingly adopted Machine Learning and LLMs as integrated tools for cancer diagnosis. In particular, studies show that Machine Learning models have high performance sometimes near to the 100\% accuracy in cancer detection\cite{Tandon2024}, with the most studied cancers being breast cancer, lung cancer, followed by brain, cervical and liver cancer, using mainly CNN as the primary architecture\cite{Tandon2024}. 
\\\\
A systematic literature review of February 2026 shows that ML approaches remain concentrated on well-studied cancer types such as breast and lung cancer, while other rare or low reproducibily cancers remain underrepresented or understudied in the literature\cite{rosca2026}. 
\\\\
While ML classifiers achieved reasonable success in specific tasks such as thyroid carcinoma diagnosis, their dependency on manual feature extraction limited their generalizability and robustness against data variability, creating a bottleneck in this research field. This gap led to the research in Large Language Models. Recently, their research in the oncology field has increased significantly, due to their wide applicability such as the ability to process unstructured data, their generalization capabilities and their ability to generate natural language, enabling automated extraction of relevant information from clinical texts and supporting clinical decision-making processes\cite{rahman2026}. 
This technology could also resolve a common problem of ML-based diagnostic tools, which is the lack of explainability and then, the trust of such tools by medical professionals\cite{rahman2026}. However, studies show some critical issue about the use of LLM in this field. The main issue refers to the hallucination that may produce during the generation of the responses\cite{rahman2026}.}

\paragraph{Idea}
{The idea of this project is to fill this gap by developing technologies that can be integrated in the medical workflow, with a particular focus on oncology. 
The core of the project is the design and implementation of a framework architecture that includes:
\begin{itemize}
    \item \textbf{Automation of Data Extraction}: The first important component of the architecture is the automation of data extraction from clinical documents. This component will be based on a Language Model that will be used to extract relevant information. Its usage will be focused on the extraction, but also on the structuring of such data, aiming to organize this information in well known standard formats such as ICD-O-3\footnote{\url{https://www.who.int/standards/classifications/other-classifications/international-classification-of-diseases-for-oncology}} and TNM\footnote{\url{https://www.uicc.org/resources/tnm}}.
    \item \textbf{Data Quality and Dyscrasias Detection}: The second component of the architecture includes the identification of possible human errors in the data. This component will ensure the quality of the structured data and it will be based on a Language Model that will be used to identify anomalies in already existing data. The usage of LLMs/SLMs improve the speed of identification process and it generates natural language explanations of the possible errors, facilitating the human-in-the-loop approach for the supervision of the process. 
    \item \textbf{Diagnostic Support}: If the previous components are successfully implemented and evaluated in a modest time, the architecture could be extended to include a diagnostic support component.
If the architecture has a good performance in the previous components, the logical next step is to use the produced data to train a Machine Learning model for diagnostic purpose and then integrate it in the framework.
   \item \textbf{Integration of the architecture in a real medical workflow}: A core focus of this project is ensuring high potential for real-world clinical adoption. In order to ensure that, standard formats will be used and the framework must remain generic and extendible.
In case of collaborations with public istitutes, the architecture will be adapted in order to maximize the usability for specific requests.
\end{itemize}}

\paragraph{Data Retrieval}{One of the main challenges of medical research is how the data is retrieved and manipulated. 
In particular, the data should represent real cases and be of high quality, but at the same time it should be anonymized to protect the privacy of patients.   
\\\\
A well known strategy to overcome the data retrieval problem is to collaborate with real sanitary public institutes. Such collaborations require time and includes the process of access and anonymization of the data that should be taken into consideration. For this reason, part of the project will be dedicated to the activity of searching for collaborations with these istitutes such as Ausl Romagna and IRST of Meldola and the anonymization process.
\\\\
A fallback strategy will be also defined in case the collaborations are not found in time. This strategy includes the use of publicly available datasets such as TCGA dataset\footnote{\url{https://github.com/jkefeli/tcga-path-reports}}, NCI Registry\footnote{\url{https://portal.gdc.cancer.gov/}} or Belgium Cancer Registry\footnote{\url{https://kankerregister.org/en}} (subject to a formal request) and it could also include the generation of synthetic data and the verification of its realism.
}

\paragraph{Project timeline}{The project is designed to be divided into three main parts:
\begin{enumerate}
    \item \textbf{Knowledge acquisition and collaboration search} (first year): The main focus will be the study of the litelature, following PhD courses and doing exams and concurrently searching for collaborations.
    \item \textbf{Design and implementation of the framework} (second year): The design and implementation will be conducted during this year. 
    \item \textbf{Evaluation and dissemination} (third year): The proposed framework will be evaluated and disseminated. The final thesis will be written. 
\end{enumerate}}

\bibliographystyle{unsrt}
\bibliography{bibliography}

\end{document}
